\documentclass[11pt,a4paper]{article}
\usepackage[margin=25mm,headheight=14pt]{geometry}
\usepackage{fontspec}
\setmainfont{lmroman10-regular.otf}[BoldFont=lmroman10-bold.otf,ItalicFont=lmroman10-italic.otf,BoldItalicFont=lmroman10-bolditalic.otf]
\setsansfont{lmsans10-regular.otf}
\setmonofont{lmmono10-regular.otf}[Scale=MatchLowercase]
\newfontfamily\authorsymbol{Noto Sans Symbols}[Scale=MatchUppercase]
\usepackage{amsmath,amssymb,amsthm,mathtools,booktabs}
\usepackage{microtype}
\usepackage[hidelinks,unicode]{hyperref}
\usepackage{fancyhdr}
\pagestyle{fancy}\fancyhf{}\fancyhead[L]{\small\itshape Exact amplitude elimination}\fancyhead[R]{\small\thepage}\renewcommand{\headrulewidth}{0pt}
\setlength{\parindent}{1.2em}\setlength{\parskip}{0pt}\setlength{\emergencystretch}{2em}
\newtheorem{theorem}{Theorem}\newtheorem{proposition}[theorem]{Proposition}\newtheorem{lemma}[theorem]{Lemma}\newtheorem{corollary}[theorem]{Corollary}\theoremstyle{definition}\newtheorem{definition}[theorem]{Definition}\newtheorem{example}[theorem]{Example}
\DeclareMathOperator{\rank}{rank}\DeclareMathOperator{\spanop}{span}\DeclareMathOperator{\diag}{diag}
\newcommand{\R}{\mathbb R}\newcommand{\C}{\mathbb C}\newcommand{\Z}{\mathbb Z}\newcommand{\N}{\mathbb N}\newcommand{\ip}[2]{\langle #1,#2\rangle}\newcommand{\norm}[1]{\lVert #1\rVert}\newcommand{\one}{\mathbf 1}
\hypersetup{pdftitle={Exact Amplitude Elimination under Robust and Quantized Constraints},pdfauthor={lmtlssss🜐},pdfsubject={Mathematics: optimization and homogeneous systems},pdfkeywords={homogeneous optimization, amplitude elimination, robust constraints, quantization, exact arithmetic}}
\title{\LARGE Exact Amplitude Elimination under\\Robust and Quantized Constraints}
\author{lmtlssss{\authorsymbol 🜐}}
\date{2 November 2026\\[3pt]\small\href{https://doi.org/10.5281/zenodo.23897131}{10.5281/zenodo.23897131}}
\begin{document}\maketitle\thispagestyle{plain}

\begin{abstract}
We give an exact elimination of a positive amplitude parameter from a family
of homogeneous constrained optimization problems. A design determines positive
lower and upper response coefficients and a positive homogeneous cost.
Finite robust scenarios reduce to one admissible amplitude interval. On an
affine amplitude lattice, the least feasible grid point gives the exact
reduced objective. We prove a sharp rounding-factor bound, characterize the
feasibility condition needed to transfer a continuous minimizer to a grid,
and give exact examples where quantization reverses the preferred design.
Relative coefficient bounds yield interval certificates for robust design
selection. Homogeneous linear recurrences supply a constructive source of
the response coefficients. All worked comparisons are verified with rational
arithmetic and independent finite enumeration.
\end{abstract}
\noindent\textit{Keywords.} Homogeneous optimization; parameter elimination;
robust constraints; lattice quantization; fractional objectives; exact arithmetic.

\section{The scale and the design}
Let $\mathcal P$ be an arbitrary nonempty design set and let $a>0$ be an
amplitude. For every $p\in\mathcal P$, consider positive response coefficients
$m_s(p),M_s(p)$ and positive cost coefficients $e_s(p)$, indexed by a finite
nonempty scenario set $\mathcal S$. Let $F_s,H_s>0$ be required lower and
upper limits. The constraints are
\begin{equation}\label{eq:constraints}
 a m_s(p)\ge F_s,\qquad a M_s(p)\le H_s
 \quad(s\in\mathcal S).
\end{equation}
Let the objective be
\begin{equation}\label{eq:cost}
 C(a,p)=a^\nu e(p),\qquad
 e(p)=\max_{s\in\mathcal S}e_s(p),\qquad \nu>0.
\end{equation}
The maximum describes a robust cost; a positive weighted sum of coefficients
has the same amplitude law. No convexity or continuity of $\mathcal P$ is
needed for the pointwise reduction below.

Define
\begin{equation}\label{eq:interval}
 L(p)=\max_s\frac{F_s}{m_s(p)},\qquad
 U(p)=\min_s\frac{H_s}{M_s(p)}.
\end{equation}
All coefficients are assumed finite. In particular, $L(p)>0$.

\begin{theorem}[Continuous amplitude elimination]\label{thm:continuous}
A design $p$ is feasible exactly when $L(p)\le U(p)$.
For a feasible design, the unique amplitude minimizing~\eqref{eq:cost} is
$a^*(p)=L(p)$, and the reduced value is
\begin{equation}\label{eq:reduced}
 \boxed{J(p)=e(p)L(p)^\nu.}
\end{equation}
Consequently,
\[
 \inf_{p,a\text{ satisfying }\eqref{eq:constraints}}C(a,p)
 =\inf_{p:L(p)\le U(p)}J(p).
\]
Whenever the right-hand infimum is attained, its minimizer together with
$a=L(p)$ is a minimizer of the original problem.
\end{theorem}
\begin{proof}
The lower inequalities in~\eqref{eq:constraints} are jointly equivalent to
$a\ge L(p)$; the upper inequalities are equivalent to $a\le U(p)$.
Their intersection is the closed interval $[L(p),U(p)]$, or the empty set.
For fixed $p$, the function $a\mapsto e(p)a^\nu$ is strictly increasing on
positive amplitudes. Its minimum on a nonempty interval is therefore at
$L(p)$, with value~\eqref{eq:reduced}. Taking the infimum over designs proves
the remaining statements.
\end{proof}

For a single lower constraint, a common floor $F$ and a quadratic cost,
\begin{equation}\label{eq:quotient}
 a^*(p)=\frac{F}{m(p)},\qquad
 J(p)=F^2\frac{e(p)}{m(p)^2},
\end{equation}
provided the upper bound accepts that amplitude. Thus the quotient
$e/m^2$ is the exact design objective obtained by eliminating scale.

\section{Affine amplitude lattices}
Fix $a_0\ge0$ and $\Delta>0$, and restrict amplitudes to
\[
 \mathcal A=\{a_0+k\Delta:k\in\Z_{\ge0}\}\cap\R_{>0}.
\]
Set $L_0(p)=\max\{L(p),a_0\}$ and define
\begin{equation}\label{eq:grid}
 k_\Delta(p)=\max\left\{0,
    \left\lceil\frac{L(p)-a_0}{\Delta}\right\rceil\right\},\qquad
 a_\Delta(p)=a_0+\Delta k_\Delta(p).
\end{equation}
Because $L(p)>0$, this chosen amplitude is positive even when $a_0=0$.

\begin{theorem}[Exact grid reduction]\label{thm:grid}
The grid-constrained problem at a design $p$ is feasible exactly when
$a_\Delta(p)\le U(p)$. In that case its unique minimum is
\begin{equation}\label{eq:gridobj}
 \boxed{J_\Delta(p)=e(p)a_\Delta(p)^\nu.}
\end{equation}
Moreover,
\begin{equation}\label{eq:round}
 L_0(p)\le a_\Delta(p)<L_0(p)+\Delta,
\end{equation}
so
\begin{equation}\label{eq:factor}
 1\le\frac{J_\Delta(p)}{e(p)L_0(p)^\nu}
 <\left(1+\frac{\Delta}{L_0(p)}\right)^\nu.
\end{equation}
The upper factor is a sharp supremum when the location of the interval
relative to the grid varies.
\end{theorem}
\begin{proof}
The integer $k_\Delta(p)$ is the least permitted index whose amplitude is
at least $L(p)$. Every other feasible grid point is larger and has larger
cost. Thus feasibility is exactly its acceptance by the upper bound, and
\eqref{eq:gridobj} follows. The definition of the ceiling proves
\eqref{eq:round}; positivity gives~\eqref{eq:factor}.
For sharpness, place $L(p)$ immediately above a grid point and keep the next
point below $U(p)$. The rounding increment approaches $\Delta$ from below,
so the ratio approaches the displayed upper factor.
\end{proof}

\begin{corollary}[Transfer of a continuous minimizer]\label{cor:transfer}
Suppose $p_*$ minimizes $e(p)L_0(p)^\nu$ over the designs with
$L_0(p)\le U(p)$, and suppose $a_\Delta(p_*)\le U(p_*)$.
Let $J^*_0$ and $J^*_\Delta$ be the corresponding continuous and grid
optimal values. Then
\[
 J^*_0\le J^*_\Delta
 <\left(1+\frac{\Delta}{L_0(p_*)}\right)^\nu J^*_0.
\]
\end{corollary}
\begin{proof}
Every grid-feasible choice is continuous-feasible, which gives the lower
bound. The rounded choice at $p_*$ is a feasible grid candidate, so its cost
bounds the grid optimum from above. Apply~\eqref{eq:factor}.
\end{proof}

The acceptance of the rounded point is part of the hypothesis. For example,
the interval $[101/100,102/100]$ is nonempty, while the lattice with origin
zero and spacing $1/10$ misses it. This exact example identifies the upper
feasibility condition required by Corollary~\ref{cor:transfer}.

\section{Quantization can change the selected design}
Take $F=1$, $\nu=2$, $a_0=0$ and $\Delta=1/10$.
For three designs, choose $(m,e)$ as in Table~\ref{tab:reversal}, and impose
an upper constraint $a m\le2$.

\begin{table}[h]
\centering
\begin{tabular}{@{}c c c c c c@{}}
\toprule
Design & $m$ & $e$ & $J=e/m^2$ & $a_\Delta$ & $J_\Delta$\\
\midrule
A & $1$ & $1$ & $1$ & $1$ & $1$\\
B & $21/20$ & $27/25$ & $48/49$ & $1$ & $27/25$\\
C & $4/5$ & $3/5$ & $15/16$ & $13/10$ & $507/500$\\
\bottomrule
\end{tabular}
\caption{An exact change of the minimizing design after amplitude quantization.}
\label{tab:reversal}
\end{table}

Design C uniquely minimizes the continuous objective, whereas A uniquely
minimizes the grid objective. All displayed amplitudes satisfy the upper
constraint. Thus optimizing the continuous quotient and rounding only its
winning design need not solve the joint grid problem. The exact reduction
in Theorem~\ref{thm:grid} instead evaluates the least feasible grid amplitude
for every candidate design.

For a finite design set of size $D$ and $S$ scenarios, constructing all
interval endpoints takes $O(DS)$ arithmetic operations. Evaluating one
exact ceiling and one reduced objective for each feasible design gives the
complete discrete search without enumerating the amplitude axis.
For rational inputs and $\nu=2$, every comparison and grid index can be
computed by integer arithmetic. Bit cost also depends on the lengths of
the rational numerators and denominators.

\section{Interval certificates for uncertain coefficients}
Suppose a single-scenario design has nominal coefficients $e,m>0$ and
relative enclosures
\[
 (1-\varepsilon)e\le\widetilde e\le(1+\varepsilon)e,\qquad
 (1-\delta)m\le\widetilde m\le(1+\delta)m,
\]
with $0\le\varepsilon<1$ and $0\le\delta<1$.

\begin{proposition}[Reduced-value enclosure]\label{prop:interval}
For a fixed positive floor $F$,
\begin{equation}\label{eq:uncertain}
 \frac{1-\varepsilon}{(1+\delta)^\nu}\,J
 \le F^\nu\frac{\widetilde e}{\widetilde m^\nu}
 \le\frac{1+\varepsilon}{(1-\delta)^\nu}\,J,
 \qquad J=F^\nu e/m^\nu.
\end{equation}
Both endpoints are attained by allowed coefficient choices.
If one design's upper endpoint is strictly below every competitor's lower
endpoint, it is the unique reduced-value minimizer throughout these
independent coefficient intervals, among designs that remain feasible.
\end{proposition}
\begin{proof}
The reduced value increases with its numerator and decreases with its
positive denominator. Substituting opposite interval endpoints proves the
two bounds and their attainability. The strict comparison of intervals
then orders every allowed realization of the design values.
\end{proof}

For a robust feasibility certificate, coefficient enclosures can be used
directly in~\eqref{eq:interval}: use the smallest allowed $m_s$ in the
lower constraint and the largest allowed $M_s$ in the upper constraint.
Independent positive interval uncertainty gives
\[
 L_{\rm rob}(p)=\max_s\frac{F_s}{m_s^-(p)},\qquad
 U_{\rm rob}(p)=\min_s\frac{H_s}{M_s^+(p)}.
\]
The continuous and grid theorems apply to this exact robust interval.

\begin{proposition}[Grid-index stability]\label{prop:gridstable}
Let the uncertain effective lower endpoint range in
$[L_-,L_+]$, with $L_->a_0$. If an integer $k\ge1$ satisfies
\[
 a_0+(k-1)\Delta<L_-\le L_+\le a_0+k\Delta,
\]
then every realization has the same least grid amplitude $a_0+k\Delta$.
It is robustly feasible when its value is at most the smallest allowed
upper endpoint. A lower endpoint interval crossing a grid boundary can
change the minimizing grid index under an arbitrarily small perturbation.
\end{proposition}
\begin{proof}
The displayed half-open grid cell is exactly the inverse image of $k$ under
$\lceil(L-a_0)/\Delta\rceil$. The upper-feasibility statement follows from
interval intersection. Values immediately on either side of a boundary
have different ceilings, proving the last statement.
\end{proof}

\section{Homogeneous recurrences and normalized responses}
The coefficient model can arise from a fixed linear recurrence whose
forcing and initial state share an amplitude. Let a design $p$ specify
matrices $A_n(p)$, input vectors $b_n(p)$ and an initial state $z_0(p)$.
Consider
\begin{equation}\label{eq:system}
 x_{n+1}=A_n(p)x_n+a b_n(p),\qquad x_0=a z_0(p).
\end{equation}
The matrices and the design sequence are fixed during the amplitude
comparison. Let $z_n$ denote the solution with $a=1$.

\begin{lemma}[Exact normalization]\label{lem:homogeneous}
For every $n$, $x_n=a z_n$. Every output linear in state and forcing
scales by $a$. Every quadratic functional of those outputs scales by $a^2$.
The same conclusion holds for piecewise-continuous linear differential
equations with amplitude-scaled initial state and forcing, including
homogeneous jumps $x^+=Jx^-+a k$.
\end{lemma}
\begin{proof}
The discrete statement follows by induction: substituting $x_n=a z_n$
into~\eqref{eq:system} yields $x_{n+1}=a(A_nz_n+b_n)=a z_{n+1}$.
Linearity and quadratic homogeneity give the output and cost laws.
For a differential equation, $a z(t)$ satisfies the scaled equation and
initial condition; uniqueness of the linear initial-value problem gives
the result between jumps. The stated jump rule preserves the identity.
\end{proof}

When normalized response extrema are positive, they give $m_s$ and $M_s$.
Lemma~\ref{lem:homogeneous} then allows one normalized trajectory per design
and scenario. Changing amplitude within this model reuses that trajectory
and applies Theorem~\ref{thm:continuous} or~\ref{thm:grid}.

A useful scalar instance has local minimum $a(1-d)$ with $0\le d<1$.
For floor $F>0$, the iteration $a_{n+1}=F+d a_n$ satisfies
\[
 a_n-a^*=d^n(a_0-a^*),\qquad a^*=F/(1-d).
\]
The fixed point is exactly the eliminated amplitude. Its derivation is
an elementary contraction calculation, and supplies a direct check for
an implementation that otherwise repeats the amplitude iteration.

\section{Verification and mathematical context}
The rational verifier compares the closed-form grid rule with direct
enumeration in 2400 finite robust instances, including 1533 feasible cases.
It checks the exact rounding inequality in every feasible instance,
verifies the design reversal in Table~\ref{tab:reversal}, retains the narrow
infeasible interval example, and checks 200 independently advanced rational
state-scaling steps. The figure contains twelve rational designs and
sixteen explicit amplitude grid points; its incidence data are included.

Eliminating a monotone scalar over an interval is elementary. The quotient
form belongs naturally to fractional programming, whose classical
parametric and transformation methods include Dinkelbach~\cite{dinkelbach}
and Schaible~\cite{schaible}. The contribution of this note is the explicit
joint treatment of robust interval construction, affine-grid feasibility,
sharp rounding control, design-order reversal and coefficient certificates,
with a complete exact implementation. The reduction leaves the design
optimization in its original mathematical class; it introduces no
convexity assumption.

\begin{thebibliography}{9}
\bibitem{dinkelbach} W. Dinkelbach, On nonlinear fractional programming,
Management Science 13(7) (1967), 492--498.
\href{https://doi.org/10.1287/mnsc.13.7.492}{doi:10.1287/mnsc.13.7.492}.
\bibitem{schaible} S. Schaible, Duality in fractional programming: a unified
approach, Operations Research 24(3) (1976), 452--461.
\href{https://doi.org/10.1287/opre.24.3.452}{doi:10.1287/opre.24.3.452}.
\end{thebibliography}
\end{document}
